% unidades/30-keigo-kenjogo.tex
\chapter{🤲 \ruby{謙譲語}{けんじょうご} — Keigo: Verbos Humildes}
\unidadheader{30}{謙譲語}{申す・いたす・参る・いただく}{Expresar humildad bajando las propias acciones}

\begin{tcolorbox}[vocabbox, title={\emoji{pushpin} Vocabulario}]
\begin{tabularx}{\linewidth}{llX}
\toprule
\textbf{Japonés} & \textbf{Español} & \textbf{Tipo} \\ \midrule
\vocabitem{申す}{もうす}{decir (humilde)}{v. G1 irr.}
\vocabitem{いたす}{いたす}{hacer (humilde)}{v. G1 irr.}
\vocabitem{参る}{まいる}{ir/venir (humilde)}{v. G1 irr.}
\vocabitem{いただく}{いただく}{recibir/comer/beber (humilde)}{v. G1}
\vocabitem{拝見する}{はいけんする}{ver (humilde)}{v. suru}
\vocabitem{存じる}{ぞんじる}{saber/conocer (humilde)}{v. G2}
\vocabitem{承る}{うけたまわる}{recibir (pedido) - muy humilde}{v. G1}
\bottomrule
\end{tabularx}
\end{tcolorbox}

\subsection*{\emoji{speech-balloon} Frases Clave}
\frase{\ruby{私}{わたし}は\ruby{山田}{やまだ}と\ruby{申}{もう}します。}{\ruby{Me llamo Yamada.} (presentación formal)}
\frase{すぐに\ruby{参}{まい}ります。}{Voy enseguida. (humilde)}
\frase{\ruby{書類}{しょるい}を\ruby{拝見}{はいけん}いたします。}{Con su permiso, leeré los documentos.}
\frase{\ruby{確認}{かくにん}させていただきます。}{Permítame confirmarlo.}

\subsection*{\emoji{gear} Gramática}
\patron{申す / いたす / 参る — verbos humildes de uso frecuente}
  {\ruby{言}{い}う → \ruby{申}{もう}す \quad する → いたす \quad \ruby{行}{い}く/\ruby{来}{く}る → \ruby{参}{まい}る}
  {\ej{\ruby{私}{わたし}は\ruby{田中}{たなか}と\ruby{申}{もう}します。}{Me llamo Tanaka.}
   \ej{ご\ruby{説明}{せつめい}いたします。}{Le explicaré. (humilde)}
   \ej{\ruby{明日}{あした}10\ruby{時}{じ}に\ruby{参}{まい}ります。}{Iré mañana a las 10. (humilde)}}

\patron{いただく / 拝見する / 存じる — recibir, ver, saber (humilde)}
  {\ruby{もらう}{もらう}/\ruby{食}{た}べる/\ruby{飲}{の}む → いただく \quad \ruby{見}{み}る → \ruby{拝見}{はいけん}する \quad \ruby{知}{し}る/\ruby{思}{おも}う → \ruby{存}{ぞん}じる}
  {\ej{\ruby{資料}{しりょう}を\ruby{拝見}{はいけん}いたしました。}{He revisado los materiales.}
   \ej{\ruby{先生}{せんせい}からご\ruby{指導}{しどう}をいただきました。}{Recibí orientación del profesor.}
   \ej{そのことは\ruby{存}{ぞん}じております。}{Tengo conocimiento de ese asunto.}}

\patron{お〜する / ご〜する — patrón humilde general}
  {お + V-stem + する \quad / \quad ご + N (kango) + する}
  {\ej{\ruby{荷物}{にもつ}をお\ruby{持}{も}ちします。}{Le llevaré el equipaje. (お持ちする)}
   \ej{\ruby{後}{のち}ほどご\ruby{連絡}{れんらく}いたします。}{Me comunicaré con usted más tarde.}
   \ej{ご\ruby{案内}{あんない}いたします。}{Le guiaré.}}

\comparacion{いただく vs もらう / いたす vs する}
  {謙譲語\\\small situación formal con superior\\\small \ruby{受}{う}け\ruby{取}{と}りいたします\\\textit{(yo, humildemente, lo recibo)}}
  {普通語\\\small situación neutra o entre iguales\\\small \ruby{受}{う}け\ruby{取}{と}ります\\\textit{(recibo — neutral)}}
  {Las formas 謙譲語 solo se usan cuando el hablante quiere bajar su propia posición frente a alguien de mayor estatus. En conversación entre iguales (amigos, colegas) el uso de 謙譲語 suena excesivamente formal o irónicamente distante. La clave: usar en contexto de trabajo, servicio al cliente y situaciones formales.}

\coloquial{〜させていただく — la forma humilde más versátil}{
  En japonés moderno, 〜させていただく (V-causativo + いただく, «me permito hacer X») es la expresión humilde más usada en entornos corporativos: 「確認させていただきます。」「担当させていただきます。」「失礼させていただきます。」 A veces se usa en exceso (させていただきすぎ), pero es muy segura y cortés.}

\culturabox{\ruby{丁寧語}{ていねいご}・\ruby{尊敬語}{そんけいご}・\ruby{謙譲語}{けんじょうご} — Los tres pilares del keigo}{
  El keigo japonés tiene tres niveles: 丁寧語 (lenguaje educado básico: ます/です), 尊敬語 (elevar al otro: いらっしゃる, おっしゃる) y 謙譲語 (humillar al propio: 申す, いたす, 参る). Un hablante competente alterna estos niveles fluidamente según el contexto social. Los errores de keigo en el trabajo pueden dañar la imagen profesional, por eso los japoneses también estudian keigo activamente.
}

\lectura{\ruby{就職活動}{しゅうしょくかつどう}の\ruby{面接}{めんせつ}\\[0.4em]
  「\ruby{本日}{ほんじつ}は\ruby{面接}{めんせつ}の\ruby{機会}{きかい}をいただき、\ruby{誠}{まこと}にありがとうございます。」「いいえ。まず、\ruby{自己紹介}{じこしょうかい}をお\ruby{願}{ねが}いします。」「はい。\ruby{私}{わたし}は\ruby{田中}{たなか}と\ruby{申}{もう}します。\ruby{御社}{おんしゃ}のことはウェブサイトで\ruby{拝見}{はいけん}し、\ruby{大変}{たいへん}\ruby{興味}{きょうみ}を\ruby{持}{も}った\ruby{次第}{しだい}でございます。」}
{Entrevista de trabajo\\[0.4em]
  «Hoy les agradezco sinceramente la oportunidad de la entrevista.» «No hay de qué. Primero, por favor haga su autopresentación.» «Sí. Me llamo Tanaka. Tuve conocimiento de su empresa a través de su sitio web y es la razón por la que me interesé mucho.»}

\begin{tcolorbox}[ejerciciobox, title={\emoji{pencil} Ejercicios}]
\begin{enumerate}
  \item Convierte a 謙譲語: 「私は山田と言います。」「明日行きます。」「資料を見ました。」
  \item ¿いただく o もらう? Contexto: le dicen a tu jefe que recibiste su email.
  \item Escribe tu autopresentación formal usando 3 formas 謙譲語.
\end{enumerate}
\textbf{Respuestas:}
1) \ruby{私}{わたし}は\ruby{山田}{やまだ}と\ruby{申}{もう}します / \ruby{明日}{あした}\ruby{参}{まい}ります / \ruby{資料}{しりょう}を\ruby{拝見}{はいけん}しました \quad
2) いただきました (\ruby{メールをいただきました}{メールをいただきました}) \quad
3) (libre)
\end{tcolorbox}

\begin{tcolorbox}[kanjibox, title={\emoji{mahjong} Kanjis}]
\begin{tabularx}{\linewidth}{cllXX}
\toprule
\textbf{Kanji} & \textbf{On} & \textbf{Kun} & \textbf{Significado} & \textbf{Vocab} \\ \midrule
\kanjirow{謙}{けん}{—}{modesto/humilde}{\ruby{謙虚}{けんきょ} / \ruby{謙譲}{けんじょう}}
\kanjirow{譲}{じょう}{\ruby{ゆず}{る}}{ceder/humillar}{\ruby{謙譲}{けんじょう} / \ruby{譲}{ゆず}る}
\kanjirow{拝}{はい}{\ruby{おが}{む}}{venerar/humilde}{\ruby{拝見}{はいけん} / \ruby{礼拝}{れいはい}}
\kanjirow{承}{しょう}{\ruby{うけたまわ}{る}}{recibir/aceptar}{\ruby{承}{うけたまわ}る / \ruby{承認}{しょうにん}}
\bottomrule
\end{tabularx}
\end{tcolorbox}
